% Part: sets-functions-relations
% Chapter: infinite
% Section: dedekind-induction

\documentclass[../../../include/open-logic-section]{subfiles}

\begin{document}

\olfileid[fa-IR]{sfr}{infinite}{induction}
\olsection[استقرای حسابی]{جبرهای ددکیند و استقرای حسابی}

اکنون نکتهٔ اساسی این است که یک جبر ددکیند---در واقع، \emph{هر} جبر
ددکیند---می‌تواند بدیلی برای اعداد طبیعی باشد. این امر از نتیجهٔ سادهٔ زیر
ناشی می‌شود:

\begin{thm}[استقرای حسابی]\ollabel{thm:dedinfiniteinduction}
فرض کنید $N, s, o$ یک جبر ددکیند را تشکیل دهند. آنگاه برای هر مجموعهٔ
$X$:
	\begin{center}
		اگر $o \in X$ و $(\forall {n} \in N \cap X){s}({n}) \in X$، {آنگاه} $N \subseteq X$.
	\end{center}
\end{thm}

\begin{proof}
بنا بر تعریف جبر ددکیند، $N = \closureofunder{s}{o}$. اکنون اگر هم
${o} \in X$ و هم $(\forall {n} \in N)(n \in X \lif
{s}({n}) \in X)$، آنگاه $N = \closureofunder{s}{o} \subseteq X$.
\end{proof}

چون استقرا مشخصهٔ اعداد طبیعی است، نکته چنین است. هر مجموعهٔ نامتناهیِ
ددکیندی که داده شود، می‌توانیم یک جبر ددکیند تشکیل دهیم و آن جبر را بدیل
خود برای اعداد طبیعی قرار دهیم.

البته \olref{thm:dedinfiniteinduction} استقرا را با عبارت‌های
\emph{مجموعه‌شناختی} صورت‌بندی می‌کند. اما به‌آسانی می‌توانیم این اصل را
با عبارت‌هایی بیان کنیم که شاید آشناتر باشند:

\begin{cor}\ollabel{natinductionschema}
فرض کنید $N, s, o$ یک جبر ددکیند را تشکیل دهند. آنگاه برای هر فرمول
$\phi(x)$ که ممکن است پارامتر داشته باشد:
\begin{center}
	اگر $\phi(o)$ و $(\forall {n} \in N)(\phi(n)\lif
	\phi({s}({n})))$، {آنگاه} $(\forall n \in N)\phi(n)$
\end{center}
\end{cor}

\begin{proof}
بگذارید $X = \Setabs{n \in N}{\phi(n)}$، و اکنون از
\olref{thm:dedinfiniteinduction} استفاده کنید.
\end{proof}

در این نتیجه از فرمولی سخن گفتیم که «پارامتر دارد». معنای تقریبی این سخن
آن است که برای هر اشیای $c_1, \ldots, c_k$ می‌توانیم با
$\phi(x, c_1, \ldots, c_k)$ کار کنیم. به‌بیان دقیق‌تر، می‌توانیم نتیجه را
بی‌ذکر «پارامترها» چنین بیان کنیم. برای هر فرمول
$\phi(x, v_1, \ldots, v_k)$ که همهٔ متغیرهای آزادش نمایش داده شده‌اند،
داریم:
	\begin{align*}
			\forall v_1 \ldots \forall v_k((&\phi(o, v_1,\ldots, v_k) \land {}\\
			&	(\forall x \in N)(\phi(x,v_1, \ldots, v_k) \lif \phi(s(x), v_1,\ldots, v_k))) \lif {}\\
			&\hspace{3em} (\forall x \in N)\phi(x, v_1,\ldots, v_k))
	\end{align*}
آشکار است که سخن‌گفتن از «داشتن پارامتر» می‌تواند خواندن مطالب را بسیار
آسان‌تر کند. (در \olref[sth][][]{part} این شیوه را بسیار به کار خواهیم
برد.)

به جبرهای ددکیند بازگردیم: هر جبر ددکیندی که داده شود، می‌توانیم توابع
معمول حسابیِ جمع، ضرب و توان‌رسانی را نیز تعریف کنیم. اما این کار
ساده نیست و فن \emph{تعریف بازگشتی} را به کار می‌گیرد. این فنی است که
بسیار دیرتر و در بافتی بسیار کلی‌تر معرفی و توجیه خواهیم کرد. (علاقه‌مندان
شاید بخواهند پس از \olref[sth][ord-arithmetic][]{chap} به این مطلب
بازگردند، یا شاید شرحی جایگزین، مانند \citealt[صص.~95--8]{Potter2004}، را
بخوانند.) اما هرگاه $N, s, o$ یک جبر ددکیند را تشکیل دهند، در نهایت
خواهیم توانست مطالب زیر را مقرر کنیم:
\begin{align*}
	{a} + {o} &= {a} & & & {a} \times {o} &= {o} & & & {a}^{o} &= s(o)\\
{a} + {s}({b}) &= {s}({a}+{b}) &&& {a} \times {s}({b}) &= ({a}\times {b}) + {a}  & & & {a}^{{s}({b})} &= {a}^{b} \times {a}
\end{align*}
و نشان دهیم که این‌ها چنان رفتار می‌کنند که امید داریم.

\end{document}
